\chapter*{General Introduction}
\addcontentsline{toc}{chapter}{General Introduction}
\markboth{General Introduction}{}

The transition from theoretical academic study to the practice of
scientific research is a decisive stage in an engineer's training. This
report documents the work carried out during a research internship
completed as part of the first year of the engineering cycle at the
Higher Institute of Applied Sciences and Technology of Sousse (ISSAT
Sousse), hosted remotely by the Maersk Mc-Kinney Moller Institute at the
University of Southern Denmark (SDU).

Deep neural classifiers now reach expert-level accuracy on many medical
imaging tasks, but their clinical adoption remains constrained by a
question they cannot themselves answer: \emph{why} was this prediction
made? The dominant family of post-hoc explanation methods --- saliency
maps and Grad-CAM among them --- answers only \emph{where} the model is
looking. It does not say \emph{what} would have to change for the
decision to be different. For a radiologist auditing a pneumonia
classifier, a heatmap over the right lower lobe is a far weaker statement
than an image showing precisely which change to that lobe would flip the
model's verdict.

\textbf{Visual counterfactual explanations} address this gap directly.
Given a query image $x$ and a target class $c$, a counterfactual
generator produces an image $x_c$ such that a classifier $f$ predicts
$f(x_c) \approx c$, while $x_c$ remains a realistic radiograph that
preserves the patient's anatomy and any non-pathological content --- ribs,
mediastinal silhouette, tubes, pacemakers and monitoring leads. The
change between $x$ and $x_c$ \emph{is} the explanation.

This internship investigated a specific research hypothesis in that
setting: that the per-step control of a guided diffusion sampler should
be \emph{learned} rather than fixed. A hand-tuned guidance schedule
applies the same editing effort to every image, which under-edits
difficult cases and over-edits --- and thereby de-realizes --- easy ones.
The system built during this internship instead frames the choice of
per-step sampling parameters as a sequential decision problem and trains
a Proximal Policy Optimization (PPO) agent to solve it, on top of a
Stable Diffusion~1.5 latent generator adapted to chest radiography with
Low-Rank Adaptation (LoRA).

The work was organised around four research questions:

\begin{itemize}[leftmargin=*]
  \item Can a reinforcement-learning agent that controls only the
        \emph{sampling procedure} --- never the diffusion weights ---
        produce valid counterfactuals on chest radiographs?
  \item Can the edits be spatially confined to clinically plausible
        regions, rather than diffusing over the whole image?
  \item Do the resulting counterfactuals encode \emph{genuine disease
        signal}, or do they merely constitute adversarial perturbations
        against the particular classifier used for guidance?
  \item How do the standard counterfactual evaluation metrics of
        Singla~\emph{et al.}~\cite{singla2022} behave under the compute
        constraints of a two-month single-GPU project?
\end{itemize}

The third question turned out to be the scientifically decisive one, and
its answer reshaped the second half of the internship. An independent
held-out validation --- re-scoring the generated counterfactuals with
classifiers that had never participated in the guidance loop --- showed
that a system reporting near-perfect validity against its own guidance
classifier transferred very poorly to independent models. That negative
result, and the redesign it motivated, form the substantive contribution
of this work and are reported in full in Chapter~\ref{ch:results}.

The report is organised into five chapters:

\begin{itemize}[leftmargin=*]
  \item \textbf{Chapter~\ref{ch:context}} presents the internship
        context, the host research institution, and the objectives, work
        packages and expected outcomes of the project.
  \item \textbf{Chapter~\ref{ch:environment}} describes the research and
        development environment: the HPC cluster and its SLURM scheduler,
        the deep-learning frameworks, the pretrained model zoo, and the
        dataset-access tooling.
  \item \textbf{Chapter~\ref{ch:tasks}} details the tasks performed ---
        the architecture of the pipeline, its module-by-module
        construction, and the iterative debugging and redesign cycles
        that the experimental evidence forced.
  \item \textbf{Chapter~\ref{ch:results}} reports the measured results,
        discusses the strengths and the substantial limitations of the
        current system, and sets out realistic perspectives for future
        work.
  \item \textbf{Chapter~\ref{ch:breast}} documents additional,
        self-initiated work extending the pipeline to a second imaging
        domain --- mammography, on the CBIS-DDSM dataset --- undertaken to
        test whether the method and its evaluation discipline generalise
        beyond chest radiography.
\end{itemize}

The RSNA pneumonia system of Chapters~\ref{ch:context}--\ref{ch:results}
remains the core deliverable and the primary subject of this report;
Chapter~\ref{ch:breast} is reported separately, and more briefly, as
additional work rather than as a second fully-validated result standing
alongside the main one.

\clearpage
